\documentclass[conference]{IEEEtran}
\IEEEoverridecommandlockouts
% The preceding line is only needed to identify funding in the first footnote. If that is unneeded, please comment it out.
%Template version as of 6/27/2024

\usepackage{cite}
\usepackage{amsmath,amssymb,amsfonts}
\usepackage{algorithmic}
\usepackage{graphicx}
\usepackage{textcomp}
\usepackage{xcolor}
\usepackage{fontspec}
\usepackage{xeCJK}
\usepackage[hidelinks]{hyperref}
\def\BibTeX{{\rm B\kern-.05em{\sc i\kern-.025em b}\kern-.08em
    T\kern-.1667em\lower.7ex\hbox{E}\kern-.125emX}}

\setmainfont{Times New Roman}
\setsansfont{Arial}
\setmonofont{Courier New}
\setCJKmainfont[BoldFont=SimHei,ItalicFont=KaiTi]{SimSun}
\setCJKsansfont[AutoFakeSlant=0.2]{SimHei}
\setCJKmonofont{FangSong}

\begin{document}

\title{运行状态约束的气象-功率融合网络用于区域风电集群预测
\thanks{This research is supported by Smart-Grid National Science and Technology Major Project (Grant No. 2025ZD0805500).}}

\author{\IEEEauthorblockN{1\textsuperscript{st} 罗傲}
\IEEEauthorblockA{\textit{School of Information Science and}\\
\textit{Technology}\\
\textit{University of Science and}\\
\textit{Technology of China}\\
Hefei 230022, China\\
aoluo@mail.ustc.edu.cn}
\and
\IEEEauthorblockN{2\textsuperscript{nd} 赵宏伟\textsuperscript{*}}
\IEEEauthorblockA{\textit{School of Information Science and}\\
\textit{Technology}\\
\textit{University of Science and}\\
\textit{Technology of China}\\
Hefei 230022, China\\
hwzhao@ustc.edu.cn}
\and
\IEEEauthorblockN{3\textsuperscript{rd} 刘磊}
\IEEEauthorblockA{\textit{School of Information Science and}\\
\textit{Technology}\\
\textit{University of Science and}\\
\textit{Technology of China}\\
Hefei 230022, China\\
liulei13@ustc.edu.cn}
\and[\hfill\break\mbox{}\hfill]
\IEEEauthorblockN{4\textsuperscript{th} 陈秋菊}
\IEEEauthorblockA{\textit{School of Information Science and}\\
\textit{Technology}\\
\textit{University of Science and}\\
\textit{Technology of China}\\
Hefei 230022, China}
\and
\IEEEauthorblockN{5\textsuperscript{th} 李斌}
\IEEEauthorblockA{\textit{School of Information Science and}\\
\textit{Technology}\\
\textit{University of Science and}\\
\textit{Technology of China}\\
Hefei 230022, China}
}

\maketitle

\begin{abstract}
区域风电集群日内出力预测是电力调度、备用容量配置和新能源消纳的重要基础。气象驱动的预测模型通常依赖未来气象预报刻画预测日内的风场演化，但同一气象形势下的实际出力还会受到预测起点前场站出力水平、短时惯性和多场站同步变化的影响；如果模型只关注未来气象，易产生与真实运行状态不一致的偏差。为此，本文提出一种运行状态约束的气象-功率融合网络。该方法利用气象预报提取气象驱动表示，将预测起点前可用的场站级实测功率编码为运行状态表示，并在解码阶段以该状态表示约束气象信息读取。吉林区域 58 个风电场站的三年数据实验表明，所提方法在日内聚合功率预测中取得 0.0779 的归一化 MAE、0.1065 的归一化 RMSE 和 0.8036 的 \(R^2\)；相对 Weather-only，MAE 和 RMSE 分别降低 6.84\% 和 6.57\%，相对融合基线均降低约 5\%。逐预测步分析显示，预测初期收益最明显，首个预测步 MAE 降低 47.94\%。结果说明，预测起点前的真实功率状态能够有效约束气象到功率的映射，提高区域风电集群预测的状态适配能力。
\end{abstract}

\begin{IEEEkeywords}
区域风电集群预测, 风电功率预测, 气象-功率融合, 运行状态约束, 多步预测
\end{IEEEkeywords}

\section{引言}

区域风电集群日内出力预测是电力调度、备用容量配置和新能源消纳的重要基础。本文以区域内多场站聚合功率为预测对象，研究其在预测日内的变化过程。未来气象预报反映预测时段内的风场演化，预测起点前的实测功率则记录场群当前运行状态，两类信息在预测中的作用不同。区域天气传播改变各场站的风能输入，进而驱动聚合功率在预测日内变化；当前出力水平、短期惯性和多场站同步变化会影响功率轨迹的起点及近端演化，限电后的出力恢复还可能使相似天气形势对应不同的实际功率。因此，模型需要同时考虑未来风场变化和预测起点处的运行状态。

预测日内可获得覆盖全时段的网格化气象预报，实测功率则仅能使用预测起点前的记录。气象预报提供后续功率变化的外部驱动信息，但不包含预测起点处的实际出力状态；该状态需要通过起点前的功率记录加以刻画。因此，近期实测功率所反映的当前出力水平、短期惯性和多场站同步变化可作为运行状态约束，参与未来气象信息的读取。本文要解决的问题是，如何在未来气象主驱动的预测中引入这一约束。

已有风电预测研究通过风速时序建模 \cite{ref1} 和网格化空间表示 \cite{ref2} 提升预测性能。面向场群的研究进一步提出自适应时空图模型 \cite{ref3} 和层级图模型 \cite{ref4}，以及模式感知图预测器 \cite{ref5} 和全局信息图编码器 \cite{ref6}。动态更新图结构 \cite{ref7} 和学习图结构 \cite{ref8} 也得到研究。其他方向包括趋势感知图切换 \cite{ref9}、多模型融合 \cite{ref10}、基于趋势的空间信息集成 \cite{ref11}，以及面向变化图输入的持续学习 \cite{ref12}。对于未来气象信息，已有工作研究多模态 Transformer 融合 \cite{ref13} 和时空注意力融合 \cite{ref14}，并分别处理 NWP 时间滞后 \cite{ref15}、风速偏移 \cite{ref16}、预测不确定性 \cite{ref_liu_bayesian_uncertainty}、多源时空融合 \cite{ref17}，以及空间相关性与 NWP 校正 \cite{ref18}。本文关注的重点是在未来气象主驱动的日内预测中，显式利用预测起点前可观测的场站功率状态。近期功率记录了当前出力水平、短期惯性和多场站同步变化，可作为运行状态约束作用于未来气象信息的读取，从而提高预测精度。

为使预测起点前的功率状态约束未来气象读取，本文将预测过程组织为气象上下文构建、预测起点状态编码和状态条件查询解码三个功能环节。气象预报被编码为气象记忆，近期场站级功率则被转换为动态上下文，用于调节各预测步查询对气象记忆的读取。

本文主要贡献如下：
\begin{itemize}
\item 从业务时间可用性出发，将近期实测功率建模为预测起点运行状态约束，并将其作用与持续驱动日内演化的未来气象明确区分。
\item 通过记忆层条件化实现该约束：多尺度场站功率被转换为动态上下文 token，直接参与各预测步的气象记忆读取。
\item 在吉林区域 58 个风电场站数据上，通过纯气象与输出层融合对比、真实/随机/零功率控制实验，以及逐预测步误差和注意力分析形成三组互补证据，验证真实起点状态的有效性及其随预测时效变化的作用。
\end{itemize}

\begin{figure*}[htbp]
\centerline{\includegraphics[width=0.78\textwidth]{../../paper_image_use/method_new/exported/Fig01.pdf}}
\caption{业务时间协议。 近期场站功率仅在预测日开始前可用，气象预报按时间关系分为历史段与目标段；本文在该时间可用性约束下进行气象-功率融合预测。}
\label{fig:time-protocol}
\end{figure*}

\begin{figure*}[htbp]
\centerline{\includegraphics[width=0.72\textwidth]{../../paper_image_use/method_new/exported/Fig02.pdf}}
\caption{整体框架。 未来气象构成主记忆，近期场站功率生成动态上下文 token，并在 memory 层面形成对气象读取过程的运行状态约束。}
\label{fig:framework}
\end{figure*}

\section{相关工作与差异}

围绕区域或场群风电预测中的时间变化和场站间空间相关性，已有研究从风速序列建模、规则网格嵌入和图结构建模等方向展开。早期深度学习方法利用 CNN-LSTM 建模风速时序特征 \cite{ref1}，也有工作将风机嵌入规则网格，使卷积网络能够处理空间分布 \cite{ref2}。面向风电场群，Wang 等在训练中自适应调整图拓扑，并结合时空图卷积、改进多头注意力和多任务输出 \cite{ref3}；Li 等将场站划分为层级子区域，并将图建模与面向极值的校正机制结合 \cite{ref4}；Liu 等识别功率波动模式，并为不同模式训练相应的时空图预测器 \cite{ref5}。Yang 等构造多种全局信息图，根据动态相关状态切换自适应图卷积模型 \cite{ref6}；后续工作通过节点嵌入动态更新场站空间相关性 \cite{ref7}，或利用图偏差注意力学习图结构，并针对负荷峰谷时段实施动态误差校正 \cite{ref8}。另有研究构造波动相关图，并依据 NWP 有效趋势切换图结构，以同时支持确定性和概率预测 \cite{ref9}；D-stacking 将时空图特征与双重多模型融合结合，用于滚动多步预测 \cite{ref10}；趋势分类和空间信息集成被用于组织日前场群预测输入 \cite{ref11}；全局感知动态时空图持续学习框架则用于适应变化的图输入并避免遗忘历史动态 \cite{ref12}。这些研究说明，多场站功率动态和空间耦合具有重要价值，但其重点主要是历史功率表示或场群图结构构造，预测起点真实功率状态应如何约束未来气象读取仍缺少显式刻画。

围绕多步风电预测中的未来气象驱动，M2TNet 通过多任务 Transformer 联合学习多模态表示和辅助任务 \cite{ref13}，M2STAN 则利用时空注意力将多模态融合扩展到多场站预测 \cite{ref14}。针对 NWP 信息的局限性，已有工作使用序列级注意力处理气象变量的时间滞后 \cite{ref15}，并针对 NWP 风速偏移设计误差容忍策略 \cite{ref16}；深度贝叶斯模型用于量化预测不确定性 \cite{ref_liu_bayesian_uncertainty}。另有研究融合多源时空特征进行场群预测 \cite{ref17}，或将空间相关性与校正后的 NWP 结合 \cite{ref18}。对于网格化气象场，3D CNN 可直接从 NWP 数据立方体中提取特征 \cite{ref19}；WindDragon 提供了自动化区域深度学习框架 \cite{ref20}；深度 CNN 编码还被用于结合保形回归森林开展区域概率预测 \cite{ref21}。本文采用 Pangu-Weather 作为未来气象来源 \cite{ref22}，并将其视为预测日内功率变化的主驱动信息。

Informer 为长序列时间序列预测引入高效注意力机制 \cite{ref23}。PatchTST 将时间序列表示为时间片段 \cite{ref24}，Crossformer 则建模多变量序列中的跨维依赖 \cite{ref25}；Non-stationary Transformers 显式处理时间序列的分布非平稳性 \cite{ref26}。在风电预测任务中，特征-时间 Transformer 用于实现可解释的多变量建模 \cite{ref_liu_feature_temporal}，TDFormer 则通过多重依赖解耦提升多步预测精度 \cite{ref_liu_tdformer}。这些方法说明，注意力机制适合处理多步预测中的长时依赖；业务化风电预测中，未来气象驱动与预测起点运行状态之间的信息分工仍需进一步显式刻画。

已有多源融合研究通常侧重气象、功率和空间信息之间的特征互补，较少关注预测起点功率状态对未来气象读取的约束作用。本文从预测起点前的场站级实测功率中生成动态上下文表示，并将其引入气象记忆读取，使各预测步的气象信息提取受到当前运行状态的约束。

\section{问题定义与业务时间对齐}

设预测日为 \(D\)，数据时间分辨率为 1 小时。为区分气象产品的起报时刻和功率预测起点，记 \(t_{\mathrm{init}}\) 为 Pangu-Weather 的气象起报时刻，即北京时间 \(D-1\) 日 08:00；记 \(t_0\) 为功率预测起点，即 \(D-1\) 日 23:00。本文以 \(L=16\) 表示功率状态窗口包含的逐小时观测数，该窗口覆盖 \(t_0-L+1\) 至 \(t_0\)，即 \(D-1\) 日 08:00--23:00；以 \(H=24\) 表示日内预测步数，第 \(h\) 个预测步对应 \(t_0+h\)，\(h=1,\ldots,H\)，因而目标窗口覆盖 \(D\) 日 00:00--23:00。

记 \(N\) 为区域内风电场站数量，预测起点前的状态上下文为 \(\mathbf{P}^{c}\in\mathbb{R}^{N\times L}\)。区域聚合功率目标记为 \(\mathbf{y}=[y_{t_0+1},\ldots,y_{t_0+H}]\in\mathbb{R}^{H}\)，模型输出为 \(\hat{\mathbf{y}}\in\mathbb{R}^{H}\)。场站邻接矩阵 \(\mathbf{A}\in\mathbb{R}^{N\times N}\) 由同一 \(L\) 步功率状态窗口估计，保留 Pearson 相关系数大于 0.7 的正相关系数作为边权。

业务时间协议如图~\ref{fig:time-protocol} 所示。由 \(t_{\mathrm{init}}\) 起报的同一次 Pangu-Weather 产品中，预报时效 0 至 \(L-1\) 对应 \(D-1\) 日 08:00--23:00，记为气象预报历史段 \(\mathbf{X}^{h}\)；预报时效 \(L\) 至 \(L+H-1\) 对应目标窗口 \(D\) 日 00:00--23:00，记为气象预报目标段 \(\mathbf{X}^{f}\)。这里的“历史段”和“目标段”仅表示气象帧的有效时刻位于目标窗口之前或之内，二者均来自 \(t_{\mathrm{init}}\) 时已经生成的同一次气象预报。两段沿时间维拼接为 \(\mathbf{X}=[\mathbf{X}^{h};\mathbf{X}^{f}]\)。因此，以 \(t_0\) 为起点生成的功率预测只读取截至 \(t_0\) 的实测功率和此前起报的气象产品，不使用目标窗口内的实测功率。

训练阶段采用逐小时滑动窗口，并保持上述输入窗口与目标窗口的相对时间对齐；验证和测试阶段采用严格日样本，按照固定业务时间协议预测 \(D\) 日出力。

\begin{figure*}[htbp]
\centerline{\includegraphics[width=0.72\textwidth]{../../paper_image_use/method_new/exported/Fig03.pdf}}
\caption{气象记忆构建。 解耦气象分支提取风场、背景场和气象变化特征，并构造可由未来查询向量读取的时空气象记忆。}
\label{fig:weather-memory}
\end{figure*}

\section{运行状态约束的气象-功率融合网络}

\subsection{整体框架}

如图~\ref{fig:framework} 所示，该网络实现三个功能环节。首先，气象编码器（Weather Encoder）将连续气象预报转换为可查询的气象上下文，承载日内演化的主要驱动信息。其次，功率状态编码器（Power-State Encoder）仅使用预测起点前可获得的场站级功率，形成随样本变化的运行状态表示。最后，查询解码器（Query Decoder）在该状态条件下读取气象上下文，并行输出全部预测步。

模型采用记忆层条件化，将编码后的运行状态映射为气象记忆隐空间中的动态上下文 token，使各预测步查询可从同一融合记忆中自适应读取预测起点状态和未来气象信息。下文依次介绍气象记忆构建、功率状态编码以及记忆层融合与查询解码结构。

\subsection{气象记忆构建}

气象记忆构建过程如图~\ref{fig:weather-memory} 所示。气象编码器输入为 \(\mathbf{X}=[\mathbf{X}^{h};\mathbf{X}^{f}]\in\mathbb{R}^{T\times C\times H_g\times W_g}\)，其中 \(T=L+H\)，\(L\) 和 \(H\) 分别为第~III 节定义的功率状态窗口长度和日内预测长度；历史段 \(\mathbf{X}^{h}\) 和目标段 \(\mathbf{X}^{f}\) 分别包含 \(L\) 帧和 \(H\) 帧。原始 \(C=17\) 个气象通道划分为 9 个风场通道和 8 个背景气象通道，并对全部通道额外计算一阶时效差分；三类信息分别通过 CNN stem 提取逐帧空间特征：
\begin{equation}
\mathbf{F}^{u}_{t}=\phi_{u}(\mathbf{X}^{u}_{t}),\quad
\mathbf{F}^{b}_{t}=\phi_{b}(\mathbf{X}^{b}_{t}),\quad
\mathbf{F}^{\Delta}_{t}=\phi_{\Delta}(\Delta\mathbf{X}_{t}),
\end{equation}
其中，\(\mathbf{X}^{u}_{t}\) 和 \(\mathbf{X}^{b}_{t}\) 分别选取第 \(t\) 帧的风场与背景气象通道。一阶时效差分定义为
\begin{equation}
\Delta\mathbf{X}_{1}=\mathbf{0},\qquad
\Delta\mathbf{X}_{t}=\mathbf{X}_{t}-\mathbf{X}_{t-1},\quad t=2,\ldots,T.
\end{equation}
算子 \(\phi_{u}\)、\(\phi_{b}\) 和 \(\phi_{\Delta}\) 是逐帧二维卷积 stem，分别编码风场驱动、背景气象状态和气象场随预报时效的变化。

三组特征通过 \(1\times1\) 模块融合，并加入时间嵌入和阶段嵌入：
\begin{equation}
\mathbf{H}_{t}=\psi([\mathbf{F}^{u}_{t};\mathbf{F}^{b}_{t};\mathbf{F}^{\Delta}_{t}])
+\mathbf{e}^{time}_{t}+\mathbf{e}^{stage}_{t}.
\end{equation}
其中，\(\psi(\cdot)\) 表示用于混合三类气象分支通道信息的 \(1\times1\) 卷积融合块；\(\mathbf{H}_{t}\) 是第 \(t\) 个气象帧的融合隐表示；\(\mathbf{e}^{time}_{t}\) 是标识气象帧在输入序列中相对位置的时间位置嵌入；\(\mathbf{e}^{stage}_{t}\) 是区分气象预报历史段与目标段的阶段嵌入。

随后将序列 \(\mathbf{H}_{1:T}\) 输入时空卷积主干 \(\mathcal{E}_{w}\)。该主干由 depthwise 3D 时空卷积、空间下采样卷积和逐帧空间混合模块组成，在保持时间长度的同时建模局部时空变化。编码后的特征先经过 \(1\times1\times1\) memory projection，再池化为固定大小的时空网格：
\begin{equation}
\mathbf{Z}^{w}=
\operatorname{Pool}\left(\operatorname{Proj}(\mathcal{E}_{w}(\mathbf{H}_{1:T}))\right)
\in\mathbb{R}^{T_m\times H_m\times W_m\times d}.
\end{equation}
其中，\(\operatorname{Proj}\) 为 \(1\times1\times1\) 卷积投影，\(\operatorname{Pool}\) 表示自适应平均池化。为保留 memory token 的时空位置，本文将时间、行和列三类可学习位置嵌入相加：
\begin{equation}
\mathbf{E}^{pos}_{\tau,i,j}=
\mathbf{p}^{time}_{\tau}+\mathbf{p}^{row}_{i}+\mathbf{p}^{col}_{j}.
\end{equation}
最终气象记忆定义为
\begin{equation}
\begin{aligned}
\mathbf{M}^{w}
&=\operatorname{Flatten}(\mathbf{Z}^{w}+\mathbf{E}^{pos})
\in\mathbb{R}^{L_w\times d},\\
L_w&=T_mH_mW_m.
\end{aligned}
\end{equation}
其中，Flatten 将池化后的网格转换为 token 序列，\(L_w\) 为气象记忆 token 的数量。该记忆保留气象预报历史段、目标段和局部时空结构，使不同预测步的查询向量能够读取不同预报时效和空间区域的信息。

\subsection{功率状态编码}

功率状态编码过程如图~\ref{fig:power-state} 所示。功率状态编码器输入预测起点前可用的场站级功率 \(\mathbf{P}^{c}\in\mathbb{R}^{N\times L}\) 和场站邻接矩阵 \(\mathbf{A}\in\mathbb{R}^{N\times N}\)，其中 \(N\) 为风电场站数量，\(L\) 为第~III 节定义的功率状态窗口长度。

模型在 point、patch 和 global 三种时间尺度上编码场站级功率，分别刻画逐时刻空间状态、局部爬坡模式和上下文范围内的整体出力趋势。局部片段定义为
\begin{equation}
\label{eq:patch-start-zh}
a_j=1+(j-1)r,\qquad j=1,\ldots,J,
\end{equation}
\begin{equation}
\label{eq:patch-index-zh}
\mathcal{I}_{j}=\{a_j,\ldots,a_j+s-1\},
\end{equation}
\begin{equation}
\label{eq:patch-count-zh}
J=\left\lfloor\frac{L-s}{r}\right\rfloor+1.
\end{equation}
其中，\(s\) 是局部片段长度，\(r\) 是滑动步长，\(a_j\) 是第 \(j\) 个 patch 的起始索引，\(\mathcal{I}_j\) 是其时间索引集合，\(J\) 是有效 patch 总数。这些片段均不使用目标时段功率。三类 token 分别写为
\begin{equation}
\label{eq:point-token-zh}
\mathbf{o}^{pt}_{t}=g_{pt}(\operatorname{GCN}_{pt}(\mathbf{P}^{c}_{:,t},\mathbf{A})),
\quad t=1,\ldots,L,
\end{equation}
\begin{equation}
\label{eq:patch-token-zh}
\mathbf{o}^{pa}_{j}
=g_{pa}(\operatorname{GCN}_{pa}(\mathbf{P}^{c}_{:,\mathcal{I}_{j}},\mathbf{A})),
\quad j=1,\ldots,J,
\end{equation}
\begin{equation}
\label{eq:global-token-zh}
\mathbf{o}^{gl}
=g_{gl}(\operatorname{GCN}_{gl}(\mathbf{P}^{c},\mathbf{A})).
\end{equation}
图卷积的节点输出先沿场站和特征维展平，再分别通过线性投影 \(g_{pt}\)、\(g_{pa}\) 和 \(g_{gl}\) 映射到 \(d\) 维 token 空间。

为融合三种时间尺度的信息，功率编码器采用自底向上和自顶向下的跨尺度注意力。加入位置编码后，三组 token 分别记为 \(\mathbf{H}^{pt}\)、\(\mathbf{H}^{pa}\) 和 \(\mathbf{H}^{gl}\)。跨尺度交互首先使 patch tokens 读取 point tokens，并使 global token 读取 patch tokens；随后由 global token 调制 patch tokens，再由 patch tokens 调制 point tokens：
\begin{equation}
\mathbf{H}^{pa}_{bu}=\operatorname{CSA}_{pt\rightarrow pa}
(\mathbf{H}^{pa},\mathbf{H}^{pt}).
\end{equation}
\begin{equation}
\mathbf{H}^{gl}_{bu}=\operatorname{CSA}_{pa\rightarrow gl}
(\mathbf{H}^{gl},\mathbf{H}^{pa}_{bu}).
\end{equation}
\begin{equation}
\mathbf{H}^{pa}_{td}=\operatorname{CSA}_{gl\rightarrow pa}
(\mathbf{H}^{pa}_{bu},\mathbf{H}^{gl}_{bu}).
\end{equation}
\begin{equation}
\mathbf{H}^{pt}_{td}=\operatorname{CSA}_{pa\rightarrow pt}
(\mathbf{H}^{pt},\mathbf{H}^{pa}_{td}).
\end{equation}
其中，下标 \(bu\) 和 \(td\) 分别表示自底向上和自顶向下交互结果；\(\operatorname{CSA}_{source\rightarrow target}\) 表示以 target tokens 为 query、以 source tokens 为 key 和 value，并包含残差连接、LayerNorm 和 FFN 的跨尺度注意力块。融合后的跨尺度 token 序列为
\begin{equation}
\mathbf{S}^{p}=[\mathbf{H}^{gl}_{bu};\mathbf{H}^{pa}_{td};\mathbf{H}^{pt}_{td}]
\in\mathbb{R}^{(1+J+L)\times d}.
\end{equation}
该序列加入融合位置编码后，经过 Transformer encoder，展平并投影为样本级功率状态向量：
\begin{equation}
\mathbf{z}^{p}=f_{p}(\operatorname{vec}(\operatorname{Transformer}(\mathbf{S}^{p})))
\in\mathbb{R}^{d}.
\end{equation}
其中，\(\operatorname{vec}\) 表示展平 token 序列，\(f_p\) 为映射到 \(d\) 维的状态投影头。功率编码器由此得到样本特定的运行状态表示，将当前出力水平、局部变化和场站间耦合关系压缩后供后续气象读取过程使用。

\begin{figure*}[htbp]
\centerline{\includegraphics[width=0.66\textwidth]{../../paper_image_use/method_new/exported/Fig04.pdf}}
\caption{功率状态编码。 多尺度 GCN 从近期场站级功率中提取当前运行状态、局部变化和多场站耦合关系。}
\label{fig:power-state}
\end{figure*}

\subsection{记忆层融合与查询解码}

为在记忆层写入功率状态，两层前馈 power-to-context adapter \(a_{\theta}\) 将功率状态向量映射为 \(K\) 个动态功率上下文 token：
\begin{equation}
\mathbf{C}^{p}=\operatorname{reshape}(a_{\theta}(\mathbf{z}^{p}))\in\mathbb{R}^{K\times d}.
\end{equation}
这些 token 由同一功率状态向量联合生成，与历史功率时刻不存在逐一对应关系，随后被前置到气象记忆：
\begin{equation}
\mathbf{M}^{pc}=[\mathbf{C}^{p};\mathbf{M}^{w}]\in\mathbb{R}^{(K+L_w)\times d}.
\end{equation}

这种前置注入使功率状态以 memory token 的形式参与所有预测步的 cross-attention。不同预测步的查询向量可以在同一融合记忆中自适应读取功率上下文和气象记忆，功率状态由此作为运行状态约束参与并行解码。

查询解码器从可学习的未来查询向量 \(\mathbf{Q}^{(0)}\in\mathbb{R}^{H\times d}\) 开始，每个查询向量对应一个预测步。解码器由若干 decoder block 组成；对于 \(\ell=1,\ldots,L_d\)，每个 block 由带残差连接和 LayerNorm 的 query self-attention、query-to-memory cross-attention 和 FFN 组成。self-attention 建模不同预测步之间的依赖，cross-attention 从融合记忆 \(\mathbf{M}^{pc}\) 中读取气象记忆和功率上下文：
\begin{equation}
\label{eq:decoder-self-zh}
\tilde{\mathbf{Q}}^{(\ell)}=\operatorname{LN}\!\left(
\mathbf{Q}^{(\ell-1)}+\operatorname{MSA}(\mathbf{Q}^{(\ell-1)})\right).
\end{equation}
\begin{equation}
\label{eq:decoder-cross-zh}
\bar{\mathbf{Q}}^{(\ell)}=\operatorname{LN}\!\left(
\tilde{\mathbf{Q}}^{(\ell)}+\operatorname{MCA}(\tilde{\mathbf{Q}}^{(\ell)},\mathbf{M}^{pc})\right).
\end{equation}
\begin{equation}
\label{eq:decoder-ffn-zh}
\mathbf{Q}^{(\ell)}=\operatorname{LN}\!\left(\bar{\mathbf{Q}}^{(\ell)}
+\operatorname{FFN}(\bar{\mathbf{Q}}^{(\ell)})\right).
\end{equation}
其中，\(\operatorname{MSA}(\mathbf{Q})=\operatorname{MHA}(\mathbf{Q},\mathbf{Q},\mathbf{Q})\) 表示多头自注意力，\(\operatorname{MCA}(\mathbf{Q},\mathbf{M})=\operatorname{MHA}(\mathbf{Q},\mathbf{M},\mathbf{M})\) 表示以 \(\mathbf{Q}\) 为 query、以记忆 \(\mathbf{M}\) 为 key 和 value 的多头交叉注意力。式~\eqref{eq:decoder-self-zh}--\eqref{eq:decoder-ffn-zh} 分别执行预测步间自注意力、条件记忆读取和前馈更新；\(\tilde{\mathbf{Q}}^{(\ell)}\) 与 \(\bar{\mathbf{Q}}^{(\ell)}\) 是第 \(\ell\) 个 decoder block 中的两个中间查询表示。

最终预测为
\begin{equation}
\hat{\mathbf{y}}=g_{\phi}(\operatorname{LN}(\mathbf{Q}^{(L_d)})).
\end{equation}
其中，\(g_{\phi}\) 是预测头，将每个 horizon 的查询表示映射为对应的区域聚合功率输出，其实现由线性层、GELU、dropout 和输出线性层组成。

训练仅使用最终区域聚合功率预测损失：
\begin{equation}
\mathcal{L}=\lambda_1\,\operatorname{MAE}(\hat{\mathbf{y}},\mathbf{y})
+\lambda_2\,\operatorname{RMSE}(\hat{\mathbf{y}},\mathbf{y}).
\end{equation}
\(\lambda_1\) 和 \(\lambda_2\) 用于平衡 MAE 与 RMSE 项，本文设置 \(\lambda_1=0.6\)、\(\lambda_2=0.4\)。

所有未来查询向量均并行解码，避免自回归误差累积。训练目标直接对应最终的日内区域风电预测任务。

\section{实验}

\subsection{数据与设置}

本文基于吉林省区域风电场群构建实验数据集。功率数据覆盖 2022--2024 年，预测对象为 58 个风电场站形成的区域聚合出力；气象输入采用覆盖吉林省及周边区域的 Pangu-Weather 逐小时网格化预报，网格大小为 \(19\times33\)，空间分辨率为 \(0.25^\circ\)。17 个输入通道包括 10 m 风场、不同气压层风场、温度、湿度、位势高度和海平面气压等变量。所有功率和气象数据均按北京时间对齐。

训练集、验证集和测试集按自然年隔离，分别使用 2022、2023 和 2024 年数据；训练阶段使用滑动窗口样本，验证和测试阶段使用严格日样本。所有模型采用相同的数据划分、归一化方式和日内预测目标。目标功率按训练集预测时段功率最大值缩放，历史场站功率按场站标准化，气象输入按通道标准化；验证集和测试集仅复用训练集统计量。评价指标包括归一化 MAE、归一化 RMSE、工程尺度 MAE 和 RMSE 以及 \(R^2\)，文中结果均为 12 次独立重复实验的平均值。本文方法使用预测起点前最近 \(L\) 步场站级功率作为运行状态上下文，并由同一窗口估计邻接矩阵，输出 \(H\) 步日内聚合功率。

为检验运行状态约束的增益，本文首先设置 Weather-only 作为共享气象输入和 query decoder 的直接对比模型；同时引入 CNN、3D CNN、ViT、Swin 和 Video Transformer 等气象编码基线，用于观察不同网格气象编码方式下的纯气象预测表现。

为保证时间信息使用一致，所有 weather-only 基线均输入完整的气象预报（历史段）和气象预报（目标段）。CNN、ViT 和 Swin 基线采用逐帧空间编码后接 temporal Transformer 的方式建模时间依赖；3D CNN 和 Video Transformer 则直接进行时空编码。iTransformer-Fusion 和 PatchTST-Fusion 在 Weather-only 之外加入历史功率时序分支，并通过 softmax 门控 MLP 融合两路输出，用于比较输出层融合与本文的记忆层状态注入。

\subsection{整体预测性能}

表~\ref{tab:benchmark-short} 给出了不同模型在 2024 年独立测试集上的整体结果。本文方法取得最低归一化 MAE、最低归一化 RMSE 和最高 \(R^2\)，分别为 0.0779、0.1065 和 0.8036；对应的工程尺度 MAE 和 RMSE 分别为 383.5 MW 和 524.2 MW。Weather-only 是误差最低的纯气象模型，表明完整未来气象预报已经提供了主要预测信号。

\begin{table}[htbp]
\caption{整体预测性能}
\begin{center}
\footnotesize
\begin{tabular}{lccc}
\hline
Model & MAE \(\downarrow\) & RMSE \(\downarrow\) & \(R^2\) \(\uparrow\) \\
\hline
Ours & \textbf{0.0779} & \textbf{0.1065} & \textbf{0.8036} \\
iTrans.-Fusion & 0.0821 & 0.1121 & 0.7821 \\
PatchTST-Fusion & 0.0821 & 0.1120 & 0.7824 \\
Weather-only & 0.0836 & 0.1140 & 0.7749 \\
3D CNN & 0.0856 & 0.1172 & 0.7619 \\
CNN & 0.0916 & 0.1248 & 0.7302 \\
Video Trans. & 0.0937 & 0.1265 & 0.7228 \\
ViT & 0.0941 & 0.1278 & 0.7167 \\
Swin & 0.0947 & 0.1279 & 0.7162 \\
\hline
\end{tabular}
\label{tab:benchmark-short}
\end{center}
\end{table}

Weather-only 在纯气象模型中误差最低，两个输出层融合基线进一步利用历史功率后仍弱于本文方法。相对 Weather-only，本文方法将 MAE 和 RMSE 分别降低 6.84\% 和 6.57\%；相对 iTransformer-Fusion，两项误差分别降低 5.16\% 和 5.05\%；相对 PatchTST-Fusion 分别降低 5.19\% 和 4.99\%。结果表明，预测起点前真实功率状态在记忆读取阶段参与预测，比在输出层融合独立历史功率分支取得更低误差。

\subsection{逐预测步融合效果}

表~\ref{tab:horizon-short} 和图~\ref{fig:horizon} 展示了本文方法与 Weather-only 在不同预测步长上的误差差异，其中 \(h\) 表示第 \(h\) 个预测步。在 \(h=1\) 上，本文方法的 MAE 为 0.0445，低于 Weather-only 的 0.0856；对应 RMSE 分别为 0.0626 和 0.1138。

\begin{table}[htbp]
\caption{逐预测步误差对比}
\begin{center}
\footnotesize
\begin{tabular}{llccccc}
\hline
Model & Metric & \(h=1\) & \(h=6\) & \(h=12\) & \(h=H\) & Avg \\
\hline
Ours & MAE & \textbf{0.0445} & \textbf{0.0759} & \textbf{0.0742} & \textbf{0.0955} & \textbf{0.0779} \\
Weather & MAE & 0.0856 & 0.0806 & 0.0754 & 0.0973 & 0.0836 \\
Ours & RMSE & \textbf{0.0626} & \textbf{0.0974} & \textbf{0.1050} & \textbf{0.1215} & \textbf{0.1055} \\
Weather & RMSE & 0.1138 & 0.1066 & 0.1082 & 0.1245 & 0.1138 \\
\hline
\end{tabular}
\label{tab:horizon-short}
\end{center}
\end{table}

\begin{figure*}[t]
\centerline{\includegraphics[width=\textwidth]{figures_short/horizon_comparison_zh.pdf}}
\caption{逐预测步 MAE 及差值。(a) MAE；(b) Weather-only 与本文方法的 MAE 差值。正值表示本文方法误差更低。}
\label{fig:horizon}
\end{figure*}

逐步长曲线显示，运行状态约束的收益主要集中在预测日初始阶段，并将逐预测步平均 MAE 从 0.0836 降至 0.0779。在 \(h=1\) 上，MAE 和 RMSE 相对 Weather-only 分别降低 47.94\% 和 45.00\%；在末预测步 \(h=H\) 上仍分别降低 1.80\% 和 2.40\%。查询交叉注意力分配给功率上下文的权重占比从 \(h=1\) 的 39.56\% 降至 \(h=H\) 的 6.56\%，与近端收益更明显的误差趋势一致。表中平均 RMSE 是 \(H\) 个预测步 RMSE 的算术平均，与表~\ref{tab:benchmark-short} 在完整测试集上计算的整体 RMSE 统计口径不同。

\subsection{消融实验}

消融实验围绕四个问题组织：性能增益是否依赖预测起点前的真实功率数值，多尺度状态编码是否有效，可查询气象记忆是否必要，以及未来气象是否仍为主要驱动。Weather-only 移除功率上下文；Random power 和 Zero power 保留相同功率侧结构但破坏其中的状态信息，用于区分信息内容与额外模型容量。Basic GCN 将多尺度编码器替换为基础图卷积。No memory、No delta、Single branch、No synoptic 和 No wind 则进一步拆解气象记忆组织与气象信息组的作用。

表~\ref{tab:ablation-short} 给出了消融结果，其中 \(\Delta\) 由未四舍五入的 12 次实验均值计算。Weather-only 的 MAE 和 RMSE 分别增加 7.34\% 和 7.04\%。Random power 和 Zero power 的退化幅度与其接近，说明性能改善依赖真实运行状态的数值内容，而不是额外参数或功率通路本身。Basic GCN 的两项误差分别增加 4.88\% 和 4.74\%，说明多尺度编码能够强化状态表示，但其作用次于首先提供有效的真实状态。

\begin{table}[htbp]
\caption{关键消融结果}
\begin{center}
\footnotesize
\setlength{\tabcolsep}{2pt}
\begin{tabular}{lccccc}
\hline
Variant & MAE \(\downarrow\) & \(\Delta\)MAE & RMSE \(\downarrow\) & \(\Delta\)RMSE & \(R^2\) \(\uparrow\) \\
\hline
Ours & \textbf{0.0779} & 0.00\% & \textbf{0.1065} & 0.00\% & \textbf{0.8036} \\
Random power & 0.0835 & 7.17\% & 0.1137 & 6.78\% & 0.7760 \\
Zero power & 0.0838 & 7.61\% & 0.1141 & 7.21\% & 0.7742 \\
Basic GCN & 0.0817 & 4.88\% & 0.1115 & 4.74\% & 0.7845 \\
Weather-only & 0.0836 & 7.34\% & 0.1140 & 7.04\% & 0.7749 \\
No memory & 0.0824 & 5.86\% & 0.1128 & 5.93\% & 0.7796 \\
No delta & 0.0810 & 3.96\% & 0.1110 & 4.24\% & 0.7865 \\
Single branch & 0.0843 & 8.22\% & 0.1162 & 9.12\% & 0.7660 \\
No synoptic & 0.0858 & 10.14\% & 0.1166 & 9.50\% & 0.7645 \\
No wind & 0.1485 & 90.68\% & 0.1991 & 87.02\% & 0.3123 \\
\hline
\end{tabular}
\label{tab:ablation-short}
\end{center}
\end{table}

气象侧消融表明，No delta、No memory、Single branch 和 No synoptic 均造成稳定退化，说明气象变化、时空组织、变量角色解耦和背景天气场均有贡献。No wind 的 MAE 和 RMSE 分别增加 90.68\% 和 87.02\%，\(R^2\) 降至 0.3123，说明风场仍是主要驱动，功率上下文并未替代未来气象。

结合表~\ref{tab:benchmark-short} 的输出层融合基线，这些对照实验将主要设计选择分别隔离：真实/随机/零功率对比说明有效的起点状态信息不可缺少；输出层融合对比支持让状态约束气象记忆读取，而非仅增加一条历史功率预测分支；Basic GCN 量化了多尺度状态编码的附加收益；气象侧消融则确认条件化记忆仍由气象驱动。因此，实验支持的是本文的信息组织方式，而非将性能改善笼统归因于多个组件的累积。

\section{讨论与结论}

三年数据实验表明，在区域风电集群日内预测中，未来气象仍是主导驱动，而预测起点前真实功率状态能够提供有效的运行状态约束。本文以该状态为条件读取气象记忆，提升了日内平均预测精度，并显著降低预测初期误差。实验结果进一步验证了显式信息角色划分和记忆层条件化在协调预测起点状态与未来气象信息方面的作用。纯气象模型、输出层融合、随机/零功率控制和气象侧消融为该机制及多尺度状态编码的附加作用提供了互补证据。当前证据仍限于单一区域的确定性预测，未来将进一步研究气象预报不确定性、跨区域迁移、模型简化和极端爬坡事件。

\begin{thebibliography}{29}

\bibitem{ref1}
H. Liu, X. Mi, and Y. Li, ``Smart deep learning based wind speed prediction model using wavelet packet decomposition, convolutional neural network and convolutional long short term memory network,'' \textit{Energy Convers. Manag.}, vol. 166, pp. 120--131, 2018.

\bibitem{ref2}
R. Yu, Z. Liu, X. Li, et al., ``Scene learning: Deep convolutional networks for wind power prediction by embedding turbines into grid space,'' \textit{Appl. Energy}, vol. 238, pp. 249--257, 2019.

\bibitem{ref3}
L. Wang, J. Gao, Y. Li, and D. Wang, ``A method for ultra-short-term wind power forecasting of large-scale wind farms based on adaptive spatiotemporal graph convolution,'' \textit{Renew. Energy}, vol. 249, Art. no. 123321, 2025.

\bibitem{ref4}
M. Li, M. Yang, and Y. Yu, ``Short-term regional wind power forecast method based on hierarchical graph neural network and extreme value theory,'' \textit{Energy Convers. Manag.}, vol. 341, Art. no. 120002, 2025.

\bibitem{ref5}
X. Liu, Y. Zhang, Z. Zhen, F. Xu, F. Wang, and Z. Mi, ``Spatio-temporal graph neural network and pattern prediction based ultra-short-term power forecasting of wind farm cluster,'' \textit{IEEE Trans. Ind. Appl.}, vol. 60, no. 1, pp. 1794--1803, 2024.

\bibitem{ref6}
M. Yang, C. Ju, Y. Huang, Y. Guo, and M. Jia, ``Short-term power forecasting of wind farm cluster based on global information adaptive perceptual graph convolution network,'' \textit{IEEE Trans. Sustain. Energy}, vol. 15, no. 3, pp. 2063--2076, 2024.

\bibitem{ref7}
M. Yang, D. Wang, W. Zhang, and X. Yv, ``A centralized power prediction method for large-scale wind power clusters based on dynamic graph neural network,'' \textit{Energy}, vol. 310, Art. no. 133210, 2024.

\bibitem{ref8}
M. Yang, Y. Guo, T. Huang, et al., ``Wind farm cluster power prediction based on graph deviation attention network with learnable graph structure and dynamic error correction during load peak and valley periods,'' \textit{Energy}, vol. 312, Art. no. 133645, 2024.

\bibitem{ref9}
M. Yang, Y. Huang, Y. Guo, W. Zhang, and B. Wang, ``Ultra-short-term wind farm cluster power prediction based on FC-GCN and trend-aware switching mechanism,'' \textit{Energy}, vol. 290, Art. no. 130238, 2024.

\bibitem{ref10}
Z. Qu, J. Li, X. Hou, and J. Gui, ``A D-stacking dual-fusion, spatio-temporal graph deep neural network based on a multi-integrated overlay for short-term wind-farm cluster power multi-step prediction,'' \textit{Energy}, vol. 281, Art. no. 128289, 2023.

\bibitem{ref11}
M. Yang, Y. Jiang, C. Xu, et al., ``Day-ahead wind farm cluster power prediction based on trend categorization and spatial information integration model,'' \textit{Appl. Energy}, vol. 388, Art. no. 125580, 2025.

\bibitem{ref12}
Q. Zhang, K. Mei, Q. Zhou, et al., ``A wind farm cluster power prediction method based on globally aware dynamic spatiotemporal graph continual learning framework,'' \textit{Energy}, vol. 348, Art. no. 140591, 2026.

\bibitem{ref13}
L. Wang, Y. He, X. Liu, L. Li, and K. Shao, ``M2TNet: Multi-modal multi-task Transformer network for ultra-short-term wind power multi-step forecasting,'' \textit{Energy Rep.}, vol. 8, pp. 7628--7642, 2022.

\bibitem{ref14}
L. Wang and Y. He, ``M2STAN: Multi-modal multi-task spatiotemporal attention network for multi-location ultra-short-term wind power multi-step predictions,'' \textit{Appl. Energy}, vol. 324, Art. no. 119672, 2022.

\bibitem{ref15}
C. Liu, X. Zhang, S. Mei, Q. Zhou, and H. Fan, ``Series-wise attention network for wind power forecasting considering temporal lag of numerical weather prediction,'' \textit{Appl. Energy}, vol. 336, Art. no. 120815, 2023.

\bibitem{ref16}
M. Yang, Y. Guo, T. Huang, and W. Zhang, ``Power prediction considering NWP wind speed error tolerability: A strategy to improve the accuracy of short-term wind power prediction under wind speed offset scenarios,'' \textit{Appl. Energy}, vol. 377, Art. no. 124720, 2025.

\bibitem{ref_liu_bayesian_uncertainty}
L. Liu, J. Liu, Y. Ye, et al., ``Ultra-short-term wind power forecasting based on deep Bayesian model with uncertainty,'' \textit{Renew. Energy}, vol. 205, pp. 598--607, 2023.

\bibitem{ref17}
M. Yang, C. Han, W. Zhang, and B. Wang, ``A short-term power prediction method for wind farm cluster based on the fusion of multi-source spatiotemporal feature information,'' \textit{Energy}, vol. 294, Art. no. 130770, 2024.

\bibitem{ref18}
S. Hu, Y. Xiang, H. Zhang, et al., ``Hybrid forecasting method for wind power integrating spatial correlation and corrected numerical weather prediction,'' \textit{Appl. Energy}, vol. 293, Art. no. 116951, 2021.

\bibitem{ref19}
K. Higashiyama, Y. Fujimoto, and Y. Hayashi, ``Feature extraction of NWP data for wind power forecasting using 3D-convolutional neural networks,'' \textit{Energy Procedia}, vol. 155, pp. 350--358, 2018.

\bibitem{ref20}
J. Keisler and E. Le Naour, ``WindDragon: automated deep learning for regional wind power forecasting,'' \textit{Environ. Data Sci.}, vol. 4, Art. no. e19, 2025.

\bibitem{ref21}
J. Jonkers, D. N. Avendano, G. Van Wallendael, and S. Van Hoecke, ``A novel day-ahead regional and probabilistic wind power forecasting framework using deep CNNs and conformalized regression forests,'' \textit{Appl. Energy}, vol. 361, Art. no. 122900, 2024.

\bibitem{ref22}
K. Bi, L. Xie, H. Zhang, C. Chen, X. Gu, and Q. Tian, ``Accurate medium-range global weather forecasting with 3D neural networks,'' \textit{Nature}, vol. 619, no. 7970, pp. 533--538, 2023.

\bibitem{ref23}
H. Zhou, S. Zhang, J. Peng, et al., ``Informer: Beyond efficient Transformer for long sequence time-series forecasting,'' in \textit{Proc. AAAI Conf. Artif. Intell.}, vol. 35, no. 12, pp. 11106--11115, 2021.

\bibitem{ref24}
Y. Nie, N. H. Nguyen, P. Sinthong, and J. Kalagnanam, ``A time series is worth 64 words: Long-term forecasting with Transformers,'' in \textit{Proc. Int. Conf. Learn. Representations (ICLR)}, 2023.

\bibitem{ref25}
Y. Zhang and J. Yan, ``Crossformer: Transformer utilizing cross-dimension dependency for multivariate time series forecasting,'' in \textit{Proc. Int. Conf. Learn. Representations (ICLR)}, 2023.

\bibitem{ref26}
Y. Liu, H. Wu, J. Wang, and M. Long, ``Non-stationary Transformers: Exploring the stationarity in time series forecasting,'' in \textit{Adv. Neural Inf. Process. Syst.}, vol. 35, pp. 9881--9893, 2022.

\bibitem{ref_liu_feature_temporal}
L. Liu, X. Wang, X. Dong, et al., ``Interpretable feature-temporal transformer for short-term wind power forecasting with multivariate time series,'' \textit{Appl. Energy}, vol. 374, Art. no. 124035, 2024.

\bibitem{ref_liu_tdformer}
L. Liu, X. Liu, J. Huang, et al., ``TDFormer: A novel triple decoupled transformer for accurate multi-step wind power forecasting,'' \textit{Eng. Appl. Artif. Intell.}, vol. 181, Art. no. 115310, 2026.

\end{thebibliography}


\end{document}
